\documentclass[12pt]{article}
\usepackage[utf8]{inputenc}
\usepackage[T1]{fontenc}
\usepackage{hyperref}
\title{Yagnobi Language Resources for K-12 Curriculum Development}
\author{Compiled by Claude Sonnet 5.5 (claude-sonnet-5-5) for Andr\'as Kornai}
\date{}
\begin{document}
\maketitle

\section{Introduction}
Yagnobi (Yaghnobi; ISO 639-3 yai) is an Eastern Iranian language of Tajikistan, with about 12,000 speakers (2004) according to Wikipedia. Its UNESCO status is Definitely Endangered. Per Omniglot, it was unwritten until a Latin script was devised in the 1990s, and a Cyrillic version based on Tajik was made by Sayfiddin Mirzozoda. No school textbook was found in the searches.

\section{Primary Educational Resources}
No primers, textbooks or school programmes were found in the search results.

\section{Linguistic and Reference Materials}
\begin{itemize}
\item Wikipedia, Yaghnobi language. \url{https://en.wikipedia.org/wiki/Yaghnobi_language}
\item Omniglot, Yaghnobi, which cites A. L. Khromov, \emph{Yaghnobi Language} (Nauka, Moscow, 1972). \url{https://www.omniglot.com/writing/yaghnobi.htm}
\item A Charles University document describes a Yaghnobi--Czech dictionary expanded with fieldwork phrases and neologisms. The author is not named in the snippets. \url{https://cuni.cz/UKEN-300-version1.pdf}
\item Lubomir Novak, a Charles University Chatressar-site PDF. Its content was not verified, but it appeared in the results for this query. \url{https://chatressar.ff.cuni.cz/wp-content/uploads/sites/89/2019/01/Lubomir_Novak_26-50.pdf}
\item Numerals of Yagnobi (MPG). \url{https://lingweb.eva.mpg.de/channumerals/Yagnobi.htm}
\end{itemize}

\section{Religious and Community Resources}
None were found in the search results.

\section{Digital and Audio Resources}
Glosbe has a community-built page for the language with very limited entries. \url{https://glosbe.com/yai/dbl}

\section{Curriculum Development Context}
The searches found only reference and scholarly material, mostly in Czech or Russian. The orthography to use (Latin or Tajik Cyrillic) would have to be settled with the community. The next step would be to look for the Khromov volume, the Czech dictionary and Tajik-language materials.
\end{document}